% GENERATED FILE. Edit canonical lessons, then run npm run generate:books.
% canonical-source-hash: fnv1a64:51e0916d05c7c822
% canonical-lessons: KA-C500-svabhava, KA-C500-vyaktitva, KA-C500-nadate, KA-C500-vartane, KA-C500-anukarisu

\chapter{What Someone Is Like}
\label{ch:ka-what-someone-is-like}
\hlchaptermodality{500}

\begin{chapteropening}
\textbf{By the end of this chapter:} \emph{I can say temperament, personality, conduct, attitude and to imitate.}

Complete the last lesson of chapter 500: I can say temperament, personality, conduct, attitude and to imitate.
\end{chapteropening}

\section[\texorpdfstring{\kn{ಸ್ವಭಾವ} (svabhāva)}{ಸ್ವಭಾವ (svabhāva)}]{\kn{ಸ್ವಭಾವ} (svabhāva) — temperament; nature}

\label{lesson:KA-C500-svabhava}



\begin{quote}
Before the new one: say the Kannada for envy, then the Kannada for self-confidence.
\end{quote}

\subsection*{You'll want to know: \kn{ಸ್ವಭಾವ}}
\textbf{\kn{ಸ್ವಭಾವ}} — \emph{svabhāva} — \textquotedblleft{}temperament; nature\textquotedblright{}.

\begin{quote}
\kn{ನನ್ನ} \kn{ತಮ್ಮನ} \kn{ಸ್ವಭಾವ} \kn{ತುಂಬಾ} \kn{ಶಾಂತ}.
\end{quote}

\begin{quote}
\emph{nanna tammana svabhāva tumbā śānta.}
\end{quote}

\begin{quote}
My younger brother's nature is very calm.
\end{quote}

A person's own way of being: \kn{ಸ್ವ}, own, joined to the Sanskrit word for being.

\begin{cousinweb}
\textbf{\kn{ಸ್ವಭಾವ}} (\emph{svabhāva}) is borrowed from Sanskrit \emph{svabhāva}, \textquotedblleft{}one's own way of being, natural disposition\textquotedblright{} (Monier-Williams): \emph{sva}, \textquotedblleft{}own\textquotedblright{}, and \emph{bhāva}, \textquotedblleft{}being\textquotedblright{}.
\end{cousinweb}

\begin{grammarlens}[title={What you've built}]
One more word for describing what you lived and felt.
\end{grammarlens}

\subsection*{Your turn}
\begin{itemize}
\raggedright
  \item \emph{Say it:} \emph{svabhāva}
  \item \emph{Say it:} \emph{svabhāva}, once more
  \item \emph{Say it:} \emph{nanna tammana svabhāva tumbā śānta.} — My younger brother's nature is very calm.
  \item \emph{Recall:} say the Kannada for envy, then the Kannada for self-confidence, then say \emph{svabhāva} again
\end{itemize}

\begin{tcolorbox}[breakable,colback=teal!4,colframe=teal!35!black,title={Before you move on}]
What does \kn{ಸ್ವಭಾವ} mean? (\textquotedblleft{}Temperament; nature\textquotedblright{}.) Say it once more.
\end{tcolorbox}
\section[\texorpdfstring{\kn{ವ್ಯಕ್ತಿತ್ವ} (vyaktitva)}{ವ್ಯಕ್ತಿತ್ವ (vyaktitva)}]{\kn{ವ್ಯಕ್ತಿತ್ವ} (vyaktitva) — personality}

\label{lesson:KA-C500-vyaktitva}



\begin{quote}
Before the new one: say the Kannada for self-confidence, then the Kannada for temperament.
\end{quote}

\subsection*{You'll want to know: \kn{ವ್ಯಕ್ತಿತ್ವ}}
\textbf{\kn{ವ್ಯಕ್ತಿತ್ವ}} — \emph{vyaktitva} — \textquotedblleft{}personality\textquotedblright{}.

\begin{quote}
\kn{ಅವರ} \kn{ವ್ಯಕ್ತಿತ್ವ} \kn{ಎಲ್ಲರಿಗೂ} \kn{ಇಷ್ಟ}.
\end{quote}

\begin{quote}
\emph{avara vyaktitva ellarigū iṣṭa.}
\end{quote}

\begin{quote}
Everyone likes his personality.
\end{quote}

\kn{ವ್ಯಕ್ತಿ}, a person + -\kn{ತ್ವ}, the Sanskrit ending for \textquotedblleft{}-ness\textquotedblright{}.

\begin{cousinweb}
\textbf{\kn{ವ್ಯಕ್ತಿತ್ವ}} (\emph{vyaktitva}) is borrowed from Sanskrit \emph{vyaktitva}, \textquotedblleft{}distinctness, individuality, personality\textquotedblright{} (Monier-Williams), built on \emph{vyakti}, \textquotedblleft{}an individual\textquotedblright{}.
\end{cousinweb}

\begin{grammarlens}[title={What you've built}]
One more word for describing what you lived and felt.
\end{grammarlens}

\subsection*{Your turn}
\begin{itemize}
\raggedright
  \item \emph{Say it:} \emph{vyaktitva}
  \item \emph{Say it:} \emph{vyaktitva}, once more
  \item \emph{Say it:} \emph{avara vyaktitva ellarigū iṣṭa.} — Everyone likes his personality.
  \item \emph{Recall:} say the Kannada for self-confidence, then the Kannada for temperament, then say \emph{vyaktitva} again
\end{itemize}

\begin{tcolorbox}[breakable,colback=teal!4,colframe=teal!35!black,title={Before you move on}]
What does \kn{ವ್ಯಕ್ತಿತ್ವ} mean? (\textquotedblleft{}Personality\textquotedblright{}.) Say it once more.
\end{tcolorbox}
\section[\texorpdfstring{\kn{ನಡತೆ} (naḍate)}{ನಡತೆ (naḍate)}]{\kn{ನಡತೆ} (naḍate) — conduct; behaviour}

\label{lesson:KA-C500-nadate}



\begin{quote}
Before the new one: say the Kannada for temperament, then the Kannada for personality.
\end{quote}

\subsection*{You'll want to know: \kn{ನಡತೆ}}
\textbf{\kn{ನಡತೆ}} — \emph{naḍate} — \textquotedblleft{}conduct; behaviour\textquotedblright{}.

\begin{quote}
\kn{ಶಾಲೆಯಲ್ಲಿ} \kn{ಅವನ} \kn{ನಡತೆ} \kn{ಚೆನ್ನಾಗಿದೆ}.
\end{quote}

\begin{quote}
\emph{śāleyalli avana naḍate cennāgide.}
\end{quote}

\begin{quote}
His conduct at school is good.
\end{quote}

Built on \kn{ನಡೆ}, to walk: the way you walk through life.

\begin{cousinweb}
\textbf{\kn{ನಡತೆ}} (\emph{naḍate}) is native Dravidian: DEDR 3582 lists Kannada \emph{naḍe}, \textquotedblleft{}to walk, go, happen\textquotedblright{}, and \emph{naḍate}, \textquotedblleft{}conduct, behaviour\textquotedblright{}.
\end{cousinweb}

\begin{grammarlens}[title={What you've built}]
One more word for describing what you lived and felt.
\end{grammarlens}

\subsection*{Your turn}
\begin{itemize}
\raggedright
  \item \emph{Say it:} \emph{naḍate}
  \item \emph{Say it:} \emph{naḍate}, once more
  \item \emph{Say it:} \emph{śāleyalli avana naḍate cennāgide.} — His conduct at school is good.
  \item \emph{Recall:} say the Kannada for temperament, then the Kannada for personality, then say \emph{naḍate} again
\end{itemize}

\begin{tcolorbox}[breakable,colback=teal!4,colframe=teal!35!black,title={Before you move on}]
What does \kn{ನಡತೆ} mean? (\textquotedblleft{}Conduct; behaviour\textquotedblright{}.) Say it once more.
\end{tcolorbox}
\section[\texorpdfstring{\kn{ವರ್ತನೆ} (vartane)}{ವರ್ತನೆ (vartane)}]{\kn{ವರ್ತನೆ} (vartane) — attitude; way of behaving}

\label{lesson:KA-C500-vartane}



\begin{quote}
Before the new one: say the Kannada for personality, then the Kannada for conduct.
\end{quote}

\subsection*{You'll want to know: \kn{ವರ್ತನೆ}}
\textbf{\kn{ವರ್ತನೆ}} — \emph{vartane} — \textquotedblleft{}attitude; way of behaving\textquotedblright{}.

\begin{quote}
\kn{ಅತಿಥಿಗಳ} \kn{ಮುಂದೆ} \kn{ಮಗುವಿನ} \kn{ವರ್ತನೆ} \kn{ವಿಚಿತ್ರವಾಗಿತ್ತು}.
\end{quote}

\begin{quote}
\emph{atithigaḷa munde maguvina vartane vicitravāgittu.}
\end{quote}

\begin{quote}
The child's behaviour in front of the guests was strange.
\end{quote}

The verb is \kn{ವರ್ತಿಸು}, to behave.

\begin{cousinweb}
\textbf{\kn{ವರ್ತನೆ}} (\emph{vartane}) is borrowed from Sanskrit \emph{vartana}, \textquotedblleft{}turning, rolling\textquotedblright{} (CDIAL 11354), from the root \emph{vṛt}, \textquotedblleft{}to turn, to be\textquotedblright{}; in Kannada it is the way a person carries on.
\end{cousinweb}

\begin{grammarlens}[title={What you've built}]
One more word for describing what you lived and felt.
\end{grammarlens}

\subsection*{Your turn}
\begin{itemize}
\raggedright
  \item \emph{Say it:} \emph{vartane}
  \item \emph{Say it:} \emph{vartane}, once more
  \item \emph{Say it:} \emph{atithigaḷa munde maguvina vartane vicitravāgittu.} — The child's behaviour in front of the guests was strange.
  \item \emph{Recall:} say the Kannada for personality, then the Kannada for conduct, then say \emph{vartane} again
\end{itemize}

\begin{tcolorbox}[breakable,colback=teal!4,colframe=teal!35!black,title={Before you move on}]
What does \kn{ವರ್ತನೆ} mean? (\textquotedblleft{}Attitude; way of behaving\textquotedblright{}.) Say it once more.
\end{tcolorbox}
\section[\texorpdfstring{\kn{ಅನುಕರಿಸು} (anukarisu)}{ಅನುಕರಿಸು (anukarisu)}]{\kn{ಅನುಕರಿಸು} (anukarisu) — to imitate}

\label{lesson:KA-C500-anukarisu}



\begin{quote}
Before the new one: say the Kannada for conduct, then the Kannada for attitude.
\end{quote}

\subsection*{You'll want to know: \kn{ಅನುಕರಿಸು}}
\textbf{\kn{ಅನುಕರಿಸು}} — \emph{anukarisu} — \textquotedblleft{}to imitate\textquotedblright{}.

\begin{quote}
\kn{ಮಕ್ಕಳು} \kn{ದೊಡ್ಡವರನ್ನು} \kn{ಅನುಕರಿಸುತ್ತಾರೆ}.
\end{quote}

\begin{quote}
\emph{makkaḷu doḍḍavarannu anukarisuttāre.}
\end{quote}

\begin{quote}
Children imitate grown-ups.
\end{quote}

Past (I …): \kn{ಅನುಕರಿಸಿದೆ}, \emph{anukariside}.

The one imitated takes -\kn{ಅನ್ನು}. \kn{ಅನುಕರಣೆ} is imitation.

\begin{cousinweb}
\textbf{\kn{ಅನುಕರಿಸು}} (\emph{anukarisu}) is built on Sanskrit \emph{anukaraṇa}, \textquotedblleft{}the act of imitation\textquotedblright{} (Monier-Williams): \emph{anu}, \textquotedblleft{}after\textquotedblright{}, and \emph{kṛ}, \textquotedblleft{}to do\textquotedblright{}, with the verb ending -\emph{isu}.
\end{cousinweb}

\begin{grammarlens}[title={What you've built}]
One more word for describing what you lived and felt.
\end{grammarlens}

\subsection*{Your turn}
\begin{itemize}
\raggedright
  \item \emph{Say it:} \emph{anukarisu}
  \item \emph{Say it:} \emph{anukarisu}, once more
  \item \emph{Say it:} \emph{makkaḷu doḍḍavarannu anukarisuttāre.} — Children imitate grown-ups.
  \item \emph{Recall:} say the Kannada for conduct, then the Kannada for attitude, then say \emph{anukarisu} again
\end{itemize}

\begin{tcolorbox}[breakable,colback=teal!4,colframe=teal!35!black,title={Before you move on}]
What does \kn{ಅನುಕರಿಸು} mean? (\textquotedblleft{}To imitate\textquotedblright{}.) Say it once more.
\end{tcolorbox}
